\documentclass[12 pt, letterpaper]{hmcpset}
\usepackage{hw-preamble}
\setlist[enumerate, 1]{label = (\roman*)}

\name{Jayden Thadani}
\class{MATH 145 --- Topics in Geometry and Topology}
\assignment{Homework 7}
\duedate{11th March 2025}

\begin{document}

\begin{problem}[34.]
	Let $G = (V, E)$ be a finite graph.
	\begin{enumerate}
		\item Show that
			\[
				\sum_{x \in V} \lvert N(x) \rvert \Delta f(x) = 0
			\]
			for any scalar field $f \in \mathrm{C}^{0}(G)$.

		\item Suppose the graph is connected. Let $f \in \mathrm{C}^{0}(G)$ be a non-constant scalar field. Show that there exist points $x, y \in V$ such that $\Delta f(x) > 0$ and $\Delta f(y) < 0$.

		\item Show that (ii) fails if the graph isn't connected.
	\end{enumerate}
\end{problem}

\begin{solution}
	\begin{enumerate}
		\item Since
			 \[
				\Delta f(x) = \sum_{y \in N(x)} \frac{f(x) - f(y)}{N(x)}
			\]
			we can apply the handshake lemma to get
			\begin{align*}
				\sum_{x \in V} \lvert N(x) \rvert \Delta f(x) & = \sum_{x \in V} \sum_{y \in N(x)} f(x) - f(y) \\
									      & = \sum_{x \in V} \operatorname{deg} x \, f(x) - \sum_{y \in V} \deg y \, f(y) \\
									      & = 0.
			\end{align*}
			That is, the second equality follows because every vertex is counted as the neighbour of other vertices as many times as its degree.

		\item We showed in class that if $\Delta f = 0$ everywhere on a connected graph, then $f$ is constant. Thus, if $f$ is non-constant, there must be some point $x$ such that $\Delta f(x) > 0$ or some point $y$ such that $\Delta f(y) < 0$. Suppose without loss of generality that we have $x$ with $\Delta f(x) > 0$.

		\item Any non-constant, locally constant function on a disconnected graph is harmonic everywhere. For example, on the two point graph $G(V) = \{ 1, 2 \}$ with no edges $G(E) = \O$, each vertex has no neighbours. Thus, the sum defining the Laplacian is empty, defined to be $0$.
	\end{enumerate}
\end{solution}

\pagebreak

\begin{problem}[35.]
	Consider the Dirichlet problem for the following graph and boundary values:
	\begin{enumerate}
		\item Express the Dirichlet problem as a system of five equations in five unknowns.
		
		\item Solve the Dirichlet problem.
	\end{enumerate}
\end{problem}

\pagebreak

\begin{problem}[36.]
	Let $G = (V, E, B)$ be a finite graph-with-boundary with $\mathrm{b}_{0}(G, B) = 0$ and let $h : B \to \mathbb{R}$ be a system of boundary values. This gives rise to a Dirichlet problem. Suppose the Dirichlet problem has a \emph{symmetry}. This is a bijection $\alpha : V \to V$ with the following properties:
	\begin{itemize}
		\item It preserves the graph structure, meaning that $\alpha(x) \alpha(y) \in E$ if and only if $xy \in E$.
		\item It preserves the boundary, meaning that $\alpha(x) \in B$ if and only if $x \in B$.
		\item It preserves the boundary values, meaning that $g(\alpha(x)) = g(x)$ for all $x \in B$.
	\end{itemize}
	Let $f : V \to \mathbb{R}$ be a solution to the Dirichlet problem. Give a two-sentence proof that the symmetry preserves $f$, meaning that $f(\alpha(x)) = f(x)$ for all $x \in V$.
\end{problem}

\pagebreak

\begin{problem}[37.]
	Consider the Dirichlet problem for the following graph and boundary values:
	\begin{enumerate}
		\item Express the Dirichlet problem as a linear system of eight equations in eight unknowns.
		\item Express the Dirichlet problem as a linear system of two equations in two unknowns.

		\item Solve the Dirichlet problem.
	\end{enumerate}
\end{problem}

\pagebreak

\begin{problem}[38.]
	Consider the `infinite plane grid' graph whose vertices are the points $(m, n)$ where $m, n \in \mathbb{Z}$, and where the two vertices are connected by an edge if and only if the two points are distance $1$ apart.
	\begin{enumerate}
		\item We define six scalar fields
			\begin{align*}
				f(m, n) & = 1 & g_{1}(m, n) & = m & g_{2}(m, n) & = n \\
				h_{1}(m, n) & = m^{2} & h_{2}(m, n) & = mn & h_{3}(m, n) & = n^{2}.
			\end{align*}

			Calculate $\Delta f$, $\Delta g_{1}$, $\Delta g_{2}$, $\Delta h_{1}$, $\Delta h_{2}$, and $\Delta h_{3}$.

		\item Show that the space of harmonic scalar fields of degree at most $2$ is five-dimensional and give a basis for that space.
	\end{enumerate}
\end{problem}

\pagebreak

\begin{problem}[38.]
	Analogous to the previous question, consider the space of smooth function $\mathbb{R}^{2} \to \mathbb{R}$, and the Laplacian differential operator
	\[
		\Delta = - \left ( \frac{\partial^{2}}{\partial x^{2}} + \frac{\partial^{2}}{\partial y^{2}} \right ).
	\]
	\begin{enumerate}
		\item Show that the space of harmonic functions of degree at most $2$ is five-dimensional and give a basis for that space.

		\item Show that the $(m + 1)$-dimensional space of homogeneous polynomials of degree $m$,
			\[
				\mathcal{HP}_{m} = \{ a_{0} x^{m} + a_{1}^{m - 1} x^{m - 1} y + \dots + a_{m} y^{m} \mid a_{0}, a_{1}, \dots, a_{m} \in \mathbb{R} \}
			\]
			contains an at least $2$-dimensional subspace of functions that are harmonic.
	\end{enumerate}
\end{problem}

\begin{enumerate}
	\item The space of harmonic polynomials of degree at most $2$ (which we write as $\ker \Delta$) cannot be a full subspace of $\mathcal{P}_{2}$ (the space of all polynomials in $x$ and $y$ of degree at most $2$) since there are non-harmonic polynomials in $\mathcal{P}_{2}$
		\[
			\Delta(-x^{2}) = -1 \neq 0.
		\]
		Thus, $\ker \Delta$ is at most $5$-dimensional. We will show that it contains $5$ linearly independent harmonic polynomials, and thus, is exactly $5$-dimensional with those polynomials as a basis.

		The polynomials $1$, $x$, $y$ are clearly harmonic since they are sub-quadratic and vanish on application of any second-order differential operators. They are also linearly independent since they form a basis of the polynomials of degree at most $1$. Further, by computing
		\[
			- \Delta(xy) =  \frac{\partial^{2} (xy)}{\partial x^{2}} + \frac{\partial^{2} (xy)}{\partial y^{2}}= 0 + 0 = 0
		\]
		we see that $xy$ is also harmonic. The list $1, x, y, xy$ is still linearly independent since it is a sublist of the basis of $\mathcal{P}_{2}$. Finally,
		\[
			- \Delta(x^{2} + y^{2}) = \frac{\partial^{2} x^{2}}{\partial x^{2}} - \frac{\partial^{2} y^{2}}{\partial y^{2}} = 1 - 1 = 0.
		\]
		so $x^{2} - y^{2}$ is harmonic. The list $1, x, y, xy, x^{2} - y^{2}$ is linearly independent since the list $1, x, y, xy, x^{2}, y^{2}$ is the basis of $\mathcal{P}_{2}$.

	\item Since $\Delta$ is a differential operator of order $2$ and homogeneous (each linear term is of the same), it drops the degree of any polynomial by $2$. Thus, the image of $\mathcal{HP}_{m}$ under $\Delta$ is contained in $\mathcal{HP}_{m - 2}$. This is $2$ dimensions smaller than $\mathcal{HP}_{m}$, so by the rank-nullity theorem, $\Delta$ has a kernel of dimension at least $2$.
\end{enumerate}

\end{document}
